\section{The Verified Deception Loop}
\label{sec:design}

\RAGIN is organized around one invariant: \emph{nothing is served that has
not been either verified against the session's deception state or produced
from vetted deterministic content}. Three design decisions give the
invariant its shape, each chosen against an alternative that fails in an
instructive way.

First, verification is post-hoc. The obvious alternative is to constrain
generation itself: constrained decoding, grammar templates, or output
licenses. Such mechanisms bound the \emph{shape} of an output, but the
properties that matter here are semantic (does this banner agree with the
port from turn one?), and a promise about the process is not a
certificate about the artifact. A draft must be judged on what it says,
which is only possible after it exists; hence draft, then check.

Second, the check runs against a \emph{graph}, not against the prompt or
the transcript. \S\ref{sec:intro} argued the asymmetry: an arbitrary
answer has an unenumerable output space, but a closed world of committed
facts is small enough to check exhaustively. The fiction graph is that
closed world. Its smallness is the design's load-bearing property: it
makes verification a bounded computation with a final verdict rather
than another probabilistic judgment in need of a judge.

Third, the graph's lifetime is the session. A per-session fiction bounds
the blast radius of any drift or poisoning to one attacker's view, and an
attacker who returns presents a fresh world rather than an accumulated
liability. Cross-session continuity is a separate, host-keyed layer
described in \S\ref{sec:hostfiction}: the implementation additionally
maintains attacker-level behavioral memory in the deployed stack, and
unifying the two memory planes further
is future work (\S\ref{sec:conclusion}), not a hidden coupling.

The mechanisms below hold five goals, each naming the failure it
prevents, its enforcing component, and its falsifying measurement:
\begin{itemize}
\item \textbf{G1 (Grounding).} Deception is conditioned on who the
attacker is and on current threat intelligence---preventing irrelevant
deception no attacker class finds plausible; enforced by the profiler and
retrieval planner, measured by the paired grounding differential.
\item \textbf{G2 (Consistency).} Every observable artifact resolves to one
set of facts---preventing inconsistency as an information channel to a
re-querying attacker; enforced by the fiction graph and the verifier,
measured by contradiction rate under re-query.
\item \textbf{G3 (Safety).} No unverified content reaches the wire and
dangerous classes require human approval---preventing over-disclosure;
enforced by the verifier, canary registration, and the HITL gate,
measured by verified-serve rate and gate audit coverage.
\item \textbf{G4 (Graceful degradation).} On generation, verification, or
provider failure the loop serves deterministic vetted content---preventing
unverified improvisation under outage; enforced by the deterministic
fallback, measured by fallback correctness and serve-path provenance.
\item \textbf{G5 (Attribution).} Every decision leaves an audit trail
sufficient to reconstruct why a response was served---preventing
undebuggable behavior; enforced by serve-path provenance labels, measured
by the audit census of \S\ref{sec:evaluation}.
\end{itemize}
G5 is not a logging feature: it makes the other four falsifiable, and a
verification gate that does not persist its rejections can be believed,
not audited --- so the evaluation treats the audit trail as primary
evidence rather than instrumentation.

This section walks the loop in serving
order. Figure~\ref{fig:architecture} gives the overview; each subsection
names the goal that the component exists to satisfy.

\subsection{The Intelligence Ladder}
\label{sec:ladder}
Every mechanism in the loop is an instance of one decision: how much
intelligence to spend on the deception, and where the spending stops.
\S\ref{sec:rungs} prices that decision --- the trust price and declared
exit ramp of each rung on the certifiability axis; this subsection is
the mechanism that makes the pricing enforceable: the rungs as they
run in the loop, each with its measured effect
(\S\ref{sec:evaluation}):

\begin{itemize}
\item \textbf{R0, deterministic templates} (the fallback,
\S\ref{sec:fallback}): no intelligence, consistency by
construction, zero adaptation. This rung defines the system's floor
under failure and doubles as the control condition: ablation A5 checks that the
loop degrades to it safely, verified and serving, when the graph engine
does nothing.
\item \textbf{R1, behavioral classification} (Chrollo,
\S\ref{sec:chrollo}): low-cost routing decisions over session
features -- auditable, offline, deterministic at inference. Measured
marginal: held-out classification accuracy 0.8875 $\to$ 0.9550 with
n-gram features; a routing-accuracy ladder of its own
(\S\ref{sec:jeveval}).
\item \textbf{R2, grounded planning} (Don, \S\ref{sec:don}):
retrieval-augmented generation over curated CTI. Measured marginal: the
paired $+0.154$ composite grounding effect with 35 TTP detections
spanning 21 distinct techniques, versus zero without grounding (RQ1) -- real adaptive value, but each output remains a
statistic over an open corpus.
\item \textbf{R3, free generation under verification} (the generator
plus verifier, \S\ref{sec:verify}): the only rung whose output cannot be
certified in advance, admitted solely because the verifier converts its
artifacts into checked facts at serving time or the fallback catches
what verification rejects.
\end{itemize}

\noindent Figure~\ref{fig:ladder} renders the ladder with each rung's
measured marginal effect at its side.
The ladder runs in both directions, which makes it a design
model rather than a build history. Escalating buys capability at
a known trust price; \emph{descending} is the load-bearing direction: a
rung that can be compiled down -- replaced by a cheaper, certifiable
mechanism that reproduces its measured benefit -- should be. The routing
arm of \S\ref{sec:jeveval} shows the descent directly: an LLM
kind-assessment rung that historically fixed 44/44 misrouted sessions was
subsumed by a deterministic behavior floor at zero marginal accuracy --
the final held-out comparison scores both arms 200/200, with no
assessment output available on any session (\texttt{jev\_available}$=0$)
-- and the floor needs no network, no provider, and no trust.
Intelligence past the point
where a certifiable mechanism reproduces its effect is not an asset; it
is an attack surface with a quota bill.

% The intelligence ladder: four rungs on the certifiability axis, each
% with its measured marginal effect, and the measured fold-down from the
% top rung to the deterministic floor. Rendered from the persisted
% artifacts cited in the rung labels; all values are hardcoded here only
% because the figure must compile without the results directory.
% Geometry note: the fold-down curve is routed through the corridor left
% of the stack (control x = -3.45; every rung box edge is at x >= -2.8),
% and the stack arrows run rung-to-rung with no intervening text, so no
% arrow segment crosses a node or label interior.
% GENERATED by scripts/gen_sp_figures.py -- do not edit by hand.
% Every plotted number in the S&P figures resolves through this file, so a
% figure cannot silently disagree with the artifact it claims to show.
% Regenerate:  python scripts/gen_sp_figures.py
% Verify:      python scripts/gen_sp_figures.py --check
% Sources:
%   results/bootstrap_ci.json
%   results/stability_run1_vs_run2.json
%   results/b4_baseline_zodiac_20260922.json
%   results/model_quality/chrollo_calibration.json
%   results/sp2027_rescue/behavior_trajectory_benchmark.json

\def\RaginLiveTechnique{0.9598}
\def\RaginLiveTechniqueErrLo{0.0122}
\def\RaginLiveTechniqueErrHi{0.0122}
\def\RaginLiveActor{0.8500}
\def\RaginLiveActorErrLo{0.0000}
\def\RaginLiveActorErrHi{0.0000}
\def\RaginLivePersona{0.7325}
\def\RaginLivePersonaErrLo{0.0587}
\def\RaginLivePersonaErrHi{0.0558}
\def\RaginLiveComposite{0.8796}
\def\RaginLiveCompositeErrLo{0.0224}
\def\RaginLiveCompositeErrHi{0.0204}
\def\RaginLivePersonaRunTwo{0.7504}
\def\RaginLiveCompositeRunTwo{0.8845}
\def\RaginPersonaRunDelta{0.018}
\def\RaginTechDelta{0.317}
\def\RaginTechDeltaLo{0.162}
\def\RaginTechDeltaHi{0.480}
\def\RaginActorDelta{0.140}
\def\RaginActorDeltaLo{0.083}
\def\RaginActorDeltaHi{0.197}
\def\RaginPersonaNullDelta{0.005}
\def\RaginPersonaNullDeltaLo{-0.047}
\def\RaginPersonaNullDeltaHi{0.057}
\def\RaginOverallDelta{0.154}
\def\RaginOverallDeltaLo{0.092}
\def\RaginOverallDeltaHi{0.221}
\def\RaginGroundingTtpCount{35}
\def\RaginNoGroundingTtpCount{0}
\def\RaginRoutingBefore{0.8875}
\def\RaginRoutingAfter{0.955}

\begin{figure}[t]
\centering
\begin{tikzpicture}[
  font=\scriptsize,
  rung/.style={draw, rounded corners=1pt, thick, align=center,
               inner sep=2.5pt},
  floor/.style={rung, fill=black!12},
  gen/.style={rung, fill=black!4},
  removed/.style={rung, draw, dashed, fill=none, text=black!60},
  arr/.style={-{Latex[length=1.6mm]}, semithick},
  darr/.style={-{Latex[length=1.6mm]}, semithick, dashed}
]
  % ── rungs: mechanism on the first line, measured effect beneath it ──
  \node[removed] (r4) at (0,5.85) {kind-aware assessment
        \;{\tiny(assessed, removed)}\\
        {\tiny\itshape no marginal effect resolvable at the benchmark's ceiling}};

  \node[gen] (r3) at (0,4.05) {R3\; free generation under verification\\
        {\tiny\itshape adaptivity; 0 unverified serves across 218 queries}};

  \node[gen] (r2) at (0,2.25) {R2\; grounded planning
        \;{\tiny(Don)}\\
        {\tiny\itshape $+\RaginOverallDelta$ composite $[+\RaginOverallDeltaLo,+\RaginOverallDeltaHi]$;\; \RaginGroundingTtpCount{} vs \RaginNoGroundingTtpCount{} TTPs}};

  \node[gen] (r1) at (0,0.45) {R1\; behavioral classification
        \;{\tiny(Chrollo)}\\
        {\tiny\itshape routing $\RaginRoutingBefore \to \RaginRoutingAfter$; milliseconds, offline}};

  \node[floor] (r0) at (0,-1.35) {R0\; deterministic templates
        \;{\tiny(fallback)}\\
        {\tiny\itshape consistency by construction; zero adaptation}};

  % ── the stack: each rung rests on the cheaper one below it ──
  \draw[arr] (r0) -- (r1);
  \draw[arr] (r1) -- (r2);
  \draw[arr] (r2) -- (r3);
  \draw[arr, black!60] (r3) -- (r4);

  % ── the fold-down: measured descent from the top back to the floor,
  %    routed through the empty corridor left of the stack ──
  \draw[darr, thick]
    (r4.west) .. controls (-3.45,4.6) and (-3.45,-0.6) .. (r0.west);
  \node[rotate=90, anchor=south, font=\tiny\itshape, text=black!60,
        align=center] at (-3.62,2.0)
        {fold-down: the floor reproduces the measured benefit at zero
         trust price};
\end{tikzpicture}
\caption{The intelligence ladder. Each rung is admitted only above the
certifiable rungs beneath it, and each rung carries its measured
marginal effect (\S\ref{sec:evaluation}). The dashed arrow is the
result RQ4 adds: the assessed kind-awareness rung descended to the
deterministic floor -- no effect the benchmark resolves, no regression
risk, no provider dependency. Descending is the load-bearing direction;
a rung that cannot show a measurable benefit above the rung below it
does not stay on the ladder.}
\label{fig:ladder}
\end{figure}

\begin{figure*}[t]
\centering
% \includegraphics[width=\textwidth]{fig_architecture.pdf}
% Verified deception loop architecture (single-column TikZ)
% Routing discipline: main path is a straight vertical spine; all side
% paths run orthogonally in dedicated corridors (left margin, right gap)
% and enter nodes at distributed anchors, so no line crosses a node or
% another arrowhead.
\begin{tikzpicture}[
  node distance=3.5mm and 5mm,
  stage/.style={draw, rounded corners=1.5pt, minimum height=8mm, minimum width=21mm, align=center, font=\scriptsize},
  store/.style={stage, fill=black!8},
  gate/.style={stage, fill=black!18, thick},
  fail/.style={draw, dashed, rounded corners=1.5pt, minimum height=8mm, align=center, font=\scriptsize, fill=black!4},
  arr/.style={-{Latex[length=1.8mm]}, semithick},
  farr/.style={arr, dashed, black!60}
]
% Main serving path (top to bottom)
\node[stage] (profile) {Profile\\ \scriptsize\itshape Chrollo classifier};
\node[stage, below=of profile] (retrieve) {Retrieve + plan\\ \scriptsize\itshape Don hybrid RAG};
\node[store, below=of retrieve] (graph) {Fiction graph\\ \scriptsize\itshape facts, consistency};
\node[stage, below=of graph] (draft) {Draft response\\ \scriptsize\itshape generator LLM};
\node[gate, below=of draft] (verify) {Verify\\ \scriptsize\itshape danger + identity + consistency};
\node[stage, below=of verify] (serve) {Serve on wire};

% Side stores
\node[store, right=13mm of retrieve] (mitre) {\attck predictor\\ \scriptsize\itshape next techniques, bait};
\node[store, right=13mm of graph] (audit) {Audit log\\ \scriptsize\itshape append-only};
\node[gate, right=13mm of verify] (hitl) {HITL gate\\ \scriptsize\itshape \textsc{require} classes};

% Failure path
\node[fail, left=11mm of graph] (fallback) {Deterministic\\ fallback\\ \scriptsize\itshape vetted transcripts};

% Main path arrows (vertical spine)
\draw[arr] (profile) -- (retrieve);
\draw[arr] (retrieve) -- (graph);
\draw[arr] (graph) -- (draft);
\draw[arr] (draft) -- (verify);
\draw[arr] (verify) -- (serve) node[midway, right, font=\tiny] {pass};

% Failure arrows: orthogonal, left corridor, offset anchors
\draw[farr] (verify.west) -| ([xshift=-3mm]fallback.south)
  node[pos=0.2, above, font=\tiny] {fail $\times$2};
\draw[farr] ([xshift=3mm]fallback.south) |- (serve.west)
  node[pos=0.75, above, font=\tiny] {serve vetted};

% Approval arrow: right corridor into serve.east
\draw[arr, black!60] (hitl.south) |- (serve.east)
  node[pos=0.75, above, font=\tiny] {approved};
\draw[farr] (verify.east) -- (hitl.west) node[midway, above, font=\tiny] {\textsc{require}};

% Feedback / context arrows: right gap, offset entry into graph.south
\draw[arr, black!60] (mitre) -- (retrieve) node[midway, above, font=\tiny] {steer};
\draw[arr, black!60] (profile.east) -| (mitre.north);
\draw[arr, black!60] (graph.east) -- (audit.west) node[midway, above, font=\tiny] {record};
\draw[farr] ([xshift=8mm]draft.north) -- ([xshift=8mm]graph.south)
  node[midway, right, font=\tiny] {observe};
\end{tikzpicture}

\caption{The verified deception loop. A session enters through the
behavioral classifier (Chrollo), which routes it to a per-class deception
policy. The retrieval engine (Don) grounds the policy in CTI; the fiction
graph binds every host fact; the generator drafts under the binding; the
verifier checks the draft against the graph before serving. Dashed paths
are failure paths: verification failure falls back to deterministic
transcripts, and content flagged \textsc{require} pauses at the
human-in-the-loop gate.}
\label{fig:architecture}
\end{figure*}

\subsection{The Fiction Graph}
\label{sec:graph}
Every fact the honeypot projects about itself (hostname, kernel, listening
ports, installed databases, user accounts) is stored as a
subject--predicate--value triple in a per-session \emph{fiction graph}.
The graph is seeded from the host identity at session start and grows by
two mechanisms: \emph{declaration}, when the deception plan asserts a new
fact, and \emph{observation}, when a served response is scanned and any
host-like entities it emitted are extracted and recorded. Observation is
the mechanism that binds the LLM to the recorded state: whatever the generator said the
host is, the graph now remembers, and the next draft must agree with it.

A triple carries its source, and consistency is checked at the predicate
level: two declarations of the same predicate with different values raise
a contradiction, with normalization tuned to distinguish a genuine
conflict from a benign surface variant (a hostname is not contradicted by
its fully qualified form). When a contradiction is detected, the graph
records it and the verification node treats any draft resting on the
contradicted fact as unservable. The fiction graph is what turns G2 from
an aspiration into a data structure; the serving-time invariant below
states the consistency condition it enforces and its attacker-visible
scope.

\subsection{Verified Generation}
\label{sec:verify}
Generation is a draft, not a decision. The generator (a provider-hosted
LLM) receives the session context, the deception plan, a compact
serialization of the fiction graph (a \emph{fact block}), and the
steering context: the \attck techniques predicted to follow the
attacker's observed sequence, together with bait artifacts matched to
those techniques. The prediction uses the \attck tactic ordering to
propose the attacker's likely next moves (G1's steering half).

Before a draft reaches the wire it passes five checks:
\begin{enumerate}
\item \textbf{Danger scan.} A pattern gate rejects credential-material,
offensive-content, and real-secret signatures; one considered exemption: a
bare token a real command would print (the correct answer to
\texttt{whoami} is one word) is not chatter.
\item \textbf{Identity pin.} Identity-revealing commands answer from the
persona's pinned identity---a novice never answers \texttt{whoami} with
\texttt{root}, the hostname never changes mid-session, a root-prompting
attempt is rejected. The C1 battery carries a dedicated proof.
\item \textbf{Consistency check.} Host-like entities are observed into the
fiction graph; a contradictory draft is rejected with the contradiction
recorded.
\item \textbf{Supported-claim check.} On the declared deterministic
command surface, the draft must \emph{contain} the correct
answer---hedging, echoing, or omitting is a rejection; outside the
declared surface the check is silent, so open-ended recall is untouched.
\item \textbf{Terminal sanity.} Shell-shaped responses that answer a
command with sentence-like prose instead of executable output are
rejected.
\end{enumerate}
A draft that fails any check is retried once with the violation fed back;
a second failure triggers fallback. The checks were selected by one
criterion: each must be a \emph{decision procedure} over the draft,
decidable, cheap, and deterministic, because a verifier that itself calls
a model would only relocate the trust problem it exists to solve. What
verification deliberately does not certify is quality. A draft can pass
all five checks and still be unpersuasive; plausibility is not
decidable, and the paper does not pretend otherwise. The division of
labor is the point: the loop's machinery enforces the three properties
of the three safety constraints machine-checkably, and the consensus study of
\S\ref{sec:rq3} assesses the deception quality that only evaluation can
measure. This is the mechanism behind the zero-unverified-serves claim
evaluated in \S\ref{sec:evaluation}. The serving sequence above gives
the full loop; its fallback arm carries G4 and its \textsc{Require}
branch carries G3.

\subsection{Cross-Session Host Fiction}
\label{sec:hostfiction}
The per-session fiction buys bounded blast radius
at a measured price: a returning attacker is shown a fresh world, and
none of the session verifier's checks can see drift \emph{between}
sessions. The host-fiction layer closes exactly that gap without
widening the blast radius. It is host-keyed rather than session-keyed:
at an attacker's first contact with a host, its identity is pinned and
provisioned into every consumer --- the persona prompt is rebuilt from
the pinned identity (removing a class of contradictory-instruction
failures where prompt and verifier disagreed), the session graph is
seeded authoritatively for verification, and the fallback template
inherits the same identity. Facts observed in a session are promoted to
the host store only when they passed verification (fail-closed: an
unverified session never binds a host fact), and every subsequent
session against the same host is cross-checked against the promoted
set by a session-verifier wrapper that rejects on cross-session
contradiction. Reconnect policies are configurable: a returning
connection can share the pinned graph outright, or start a promoted
session whose facts are cross-checked against the host.

The layer is measured, not merely present. On identical deterministic
drift scenarios (reconnect drift, beyond-identity drift, cross-skill
host swap, simultaneous contradiction, repetition controls), the
session-scoped verifier alone catches 1 of 4 cross-session drift drafts;
both host-fiction policies catch 4 of 4 with 0 of 7 benign re-query
false rejections (\path{cross_session_consistency_bench.json}). The
three misses of the session-scoped baseline are invisible to it by
construction --- facts asserted in a previous session simply are not in
current-session state --- and per-skill identities name different hosts,
so a reclassified returning attacker would otherwise be shown a
different machine. On benign traffic the layer is free: paired over the
full 218-query workload, store-on versus store-off serves byte-identical
strings on every row (paired per-row delta exactly $0.0000$;
\path{external_baseline_hostfiction_comparison.json}), so cross-session
consistency is added at zero quality cost on the evaluated surface.

\subsection{Human-in-the-Loop Gate}
\label{sec:hitl}
Content classes whose compromise value is high (deployed credentials,
for instance) are marked \textsc{require} in the deception policy. When a
response would serve such content, the loop pauses and the decision is
surfaced to the operator through an audit interface that records the
context, the draft, and the decision. Every gate decision is persisted to
an append-only audit log, which \S\ref{sec:evaluation} uses as a
correctness and governance record.
The gate encodes a governance position: the system may \emph{draft}
deception, but a human decides when the deception includes trap
material (G3). The gate exists for the operator's stake, not the
attacker's: canary bait is safe to serve by construction, but it is
still an active measure whose use an operator may be required to
authorize; the audit record of that authorization leaves the system
defensible after the fact. The require-level branch is also where the
credential-deployment decisions in the benchmark trail
(Appendix~\ref{app:audit}) were made.

\subsection{Deterministic Fallback}
\label{sec:fallback}
Every scenario class has a vetted deterministic transcript: hand-written,
fiction-consistent bait and response templates (credential hunts, service
enumerations, exfiltration decoys) whose secrets are registered canaries.
When generation fails, the provider times out, verification rejects twice,
or the model refuses the persona, the loop serves the deterministic
transcript for the scenario instead. The fallback is not an admission of
defeat; it is the safety floor that makes the LLM path safe to enable.
The A5 ablation checks that the system is \emph{equivalent} through the
fallback path, and the design of the canary registration is what allows
the consensus study to distinguish deployed bait from accidental leaks.

Session state beyond identity facts replays the same discipline. A
filesystem ledger records the bytes served for every read command and
replays them verbatim when the same path is read again; time-varying
ambient content (process identifiers, timestamps, clock outputs) is not
stored but generated by hashing a per-session seed through a template,
so re-reads are identical within a session while sessions differ
plausibly from one another, and a mutating command invalidates the
affected entries so subsequent reads reflect the mutation. This ledger,
not the generative path, is what the byte-identical replay component of
the C1 result (\S\ref{sec:rq2}) exercises: determinism on re-query is a
property of how session state is stored and replayed, and the design
keeps it decoupled from generation.
The fallback also carries a hidden load: it defines the system's floor
under failure. Because the fallback serves the same fiction the generator
would have, an attacker cannot distinguish a degraded session from a
healthy one by content alone; the availability
property of G4 thereby becomes a deception property, not merely a reliability one.
What the fallback gives up is the generator's adaptivity: a fallback
session answers within its scenario's script and steers no further,
which is the correct trade only because the loop treats it as a floor
beneath the LLM path rather than the path itself.
